\subsection{Rational parametrization of the open blocks}\label{moment:open}

Write the Hermitian space as $X\oplus H\oplus X^*$ in a Witt basis,
where $X$ has ordered basis $e_1,\ldots,e_u$, $X^*$ has its dual
isotropic basis, and $H=X_0\perp Ez$ with $X_0$ standard positive
definite and $h(z,z)=-N$. Use the convention that $h$ is linear in the
second argument. Let $U$ be the unipotent radical of the stabilizer of
the complete flag in $X$. For $g$ write $D_j$ for the leading minor
of order $j$ of the matrix of pairings of $(e_1,\ldots,e_u,z)$
with its $g$-image, and $D_0=1$.

\paragraph{Rational normal form.}
If $D_j\ne0$ for $1\le j\le u$, there is a unique expression
$g=v w_p\sigma v'^{-1}$ with $v,v'\in U(E)$, $p_i\in E^\times$,
and $\sigma$ fixing the hyperbolic pairs and acting on $H$ by an
isometry. Here $w_p e_i=p_i e_i^*$, $w_p e_i^*=\bar p_i^{-1}e_i$,
and $w_p$ is the identity on $H$. Indeed, in the first hyperbolic
pair put $p=h(e_1,g e_1)$. The isotropic column $g e_1/p$ has the
form $e_1^*+y+c e_1$ in its orthogonal splitting, where
$c+\bar c=-h(y,y)$. There is a unipotent isometry $v_1$ fixing
$e_1$, sending $e_1^*$ to this column and sending a middle vector
$h_0$ to $h_0-h(y,h_0)e_1$. Hence $(v_1 w_1(p))^{-1}g$ fixes $e_1$.
Any isometry fixing $e_1$ factors as the extension of its induced
isometry on $e_1^\perp/E e_1$ times a unipotent of this same form:
after removing the induced isometry, the image of $e_1^*$ determines
$y,c$, and preservation of the pairings determines the other columns.
Iterate on the quotient. Left and right flag unipotents preserve each
leading minor (they act on the first $j$ frame vectors by
unitriangular operations, also when the last vector is $z$). Thus the
successive pivots are $p_i=D_i/D_{i-1}$ and the recursion can run to
completion. For uniqueness, two expressions have the same pivots.
The middle isometry is also determined: the relation on $X^\perp/X$
consisting of the pairs of classes of $a,g^{-1}a$ with
$a\in X^\perp\cap g X^\perp$ is invariant under the two unipotents,
and for $w_p\sigma$ it is the graph of $\sigma|_H$, with the second
coordinate as input. Finally $w_p\sigma$ conjugates $U$ to the
opposite flag unipotent; their intersection is trivial, by block
triangularity in opposite orders. This determines both unipotents.

\paragraph{The last entry and the blocks.}
For $t=h(z,\sigma z)$ write $\sigma z=\omega-(t/N)z$ with
$\omega\in X_0$. Preservation of its length gives
\[
        t\bar t-N^2=N h(\omega,\omega).
\]
The frame determinant and penultimate minor of $w_p\sigma$ are
$t\prod p_i$ and $\prod p_i$. Thus the nonzero-discriminant
condition forces $\omega\ne0$, and $|t|_v^2>N^2$ at every
archimedean place (ordinary complex absolute value). It also gives
$t=D_{u+1}/D_u$.

For a contributing cut term, both archimedean arguments have Bessel
length at most $b^{6n+2}$, with $n=u+1$. Move each independently to
its reduced box by the class congruence unipotents of
\S\ref{moment:rigidity}. Each frame column now has norm at most
$b^{37n}$. The pairings of the two frames, converted back by their
scales, have modulus at most $Nb^{74n}\le b^{134n}$ since
$N\le b^{60n}$; the corresponding rational pairings are algebraic
integers by the class-basis calculation. Their leading minors are
unchanged by the two upper unitriangular frame operations. Expansion
of a $j$-by-$j$ determinant for $j\le n\le b$ bounds its squared
modulus at every embedding by $n^{2n}b^{268n^2}\le b^{270n^2}$.
Only finitely many algebraic integers satisfy this bound. Put $A$
equal to the set of ratios of the nonzero members of this finite set. Then all
$p_i$ and $t$ lie in $A$. For fixed classes $c,c'$, reduction $s$ and
link $\Lambda$, every contributing term has the same $\sigma$ by
the graph description. Assuming absolute summability of the cut amplitudes, summing first over $U(E)\times U(E)$ and then
over the finitely many heads $p\in A^u$ gives exactly the cut block:
enforce the original class and reduction restrictions on each term.
The rational normal form proves injectivity of this parametrization
and absolute summability justifies the regrouping. Across blocks,
$(c,c',s,\sigma)$ determines $\Lambda$ and hence the block.
Consequently summing the squares of these block sums incurs no
multiplicity. It suffices to bound every finite collection and then
take the supremum.

For an integral representative in the moved frames the same elimination
works outside a finite conjugation-stable set containing the support
of the pivots, their inverses, the level, and the frame denominators:
each step divides only by its pivot or its conjugate and uses the
integral Hermitian pairing. Class changes on $X_0$ fix the hyperbolic
pairs and $z$, so the middle factor after elimination is the
class-translated middle isometry. This explains the integrality
condition on the floor-group factors in the finite-place sum. Absolute convergence and the bounds on the resulting orbital integrals enter separately.

The rational elimination and link identification above correspond to
\lineSrc{Lib/DiscreteSeries/Bruhat/BigCell}{289}{rational elimination},
\lineSrc{Lib/DiscreteSeries/Bruhat/BigCell}{374}{uniqueness}, and
\lineSrc{New/Spanning/Unfolding/Link}{309}{the graph of the middle isometry}.
The block reindexing occurs in
\lineSrc{New/Spanning/MainTerm/GlobalToLocal}{383}{the sum over heads}.

\subsection{From open-cell energy to the full moment}\label{moment:full}

The following transfer makes the dependence on the leading open-cell
error explicit. It also identifies a sufficient replacement for F6
while keeping the other errors of the manuscript. Let
$P=(\operatorname{vol}(J)\nu)^{-1/2}P_{\mathrm{raw}}$ in the
orthonormal fitted source coordinates, where $J$ is the principal
level and $\nu$ the unipotent covolume. For the unitary residue-group
action $\rho$ set
\[
 \beta(s)=\|P^*\rho(s)P\|_{\mathrm{HS}}^2,\qquad
 I_*=d_X^2/(\deg_{\mathrm{Bess}}\,m_r).
\]
Here $m_r$ is the definite genus mass. The period-to-series
identification for this $P$ says that its translated Gram blocks are
the full amplitude sums divided by $\nu$, up to the chosen convention
for complex conjugation. In particular all the sums below refer to this
single $P$. For a set $T$ of chosen places above the first level prime
and one nonzero frame vector at each place in $T$, let $K_T$ be their
simultaneous fixer and put $A_T=b^{-6k|T|}I_*$. Write
$M_T=\sum_{s\in K_T}\beta(s)$. The vector choices are included when
maximizing over $T$. Positivity of type follows from
\[
  \beta(s^{-1}t)=\langle \rho(s)PP^*\rho(s)^*,\rho(t)PP^*\rho(t)^*
                    \rangle_{\mathrm{HS}},
\]
by cyclicity of the trace.

Retain the cutoff and unequal-link error, the identity-moment bound,
the full-fixer bound and the nonopen estimate used in the original
moment argument. More precisely, at every fixer suppose the
square root of the full open mass is at most the square root of the
nonzero-discriminant link energy divided by $\nu$, plus
$b^{-15}\sqrt{I_0}$, where $I_0=|R|d_X\le A_T$ is the identity-moment scale. At the fixer with all places selected suppose $M_T\le A_T$.
At a maximizer of $b^{6k|T|}M_T$ with a place still free, assume the
nonopen mass is at most $b^{-4}M_T$. The manuscript proves the nonopen estimate by its finite-field minor
inequality, since maximality makes the mass of each additional vector
stabilizer at most $b^{-6k}M_T$. Its all-place bound uses the separate
scalar-coset and near-pair estimates. Suppose the open link energy is bounded, at every
fixer, by $(1+\delta)\nu^2 A_T$ for some $\delta\ge0$. Then
\[
    \sum_s\beta(s)\le
    \frac{(\sqrt{1+\delta}+b^{-15})^2}{1-b^{-4}} I_*.
\]
Indeed, if the maximizing fixer contains all places, comparison
with the empty fixer proves the claim directly. Otherwise the open
mass is at most $(\sqrt{1+\delta}+b^{-15})^2 A_T$ by the square-root
estimate, and $M_T$ is bounded by this plus $b^{-4}M_T$.
Divide by $1-b^{-4}$ and compare with the empty fixer.

At the constructed scales $e=(32 b x^{202})^{-1}$ satisfies
$b^{-4}=o(e)$, since $\log x=o(\log b)$. Thus
$\delta\le e/2$ suffices eventually: subtracting 1 from the displayed
factor gives $(\delta+2b^{-15}\sqrt{1+\delta}+b^{-30}+b^{-4})/
(1-b^{-4})$. In particular a uniform open-link error
$O(b^{-1-\epsilon})$ for any fixed $\epsilon>0$ suffices, provided
the other inputs just specified hold. Finally
$d_T\le |G|\deg_{\mathrm{Bess}} m_r$ converts this into
the second of the two normalized inequalities of Proposition~\ref{prop:weakerprecision}.
This propagates the relaxed target one level upstream and preserves
the required $b^{-6k|T|}$ saving at every fixer. It does not estimate
the open energy or the projected moment.

For the stronger coefficients of the pinned source, the chain is
\lineSrc{New/Spanning/Unfolding/FillingAdmissible}{743}{admissibility of the actual filling},
\lineSrc{New/Spanning/SecondMoment/SpanningTheorem}{168}{the fitted Gram blocks},
\lineSrc{New/Spanning/MainTerm/EnergyEstimate}{631}{open link energy},
\lineSrc{New/Spanning/Unfolding/SquareExpansion}{559}{the cutoff and cross-link error},
\lineSrc{New/Spanning/Unfolding/MomentBounds}{264}{open mass},
\lineSrc{New/Spanning/Unfolding/MomentBounds}{441}{maximizing fixer},
and \lineSrc{New/Spanning/SecondMoment/SpanningTheorem}{288}{target-dimension comparison}.
The source obtains $1+b^{-4}$ for link energy, $1+4b^{-4}$ for open
mass, and $1+9b^{-4}$ for full mass. The threshold for the projected
estimate must still be verified separately on $P$.

\paragraph{The outstanding estimate on the actual filling.}
For each finite collection of distinct blocks as above in one fixer,
write $B_j$ for its sum of cut double-coset amplitudes over heads.
The unresolved analytic assertion at this point is
\[
   \sum_j\|B_j\|^2
       \le (1+e/2)b^{-6k|T|} I_*\nu^2.
\]
Its quantifiers include every allowed choice of fixed vectors. The
source's stronger argument groups by $t$; the Gram matrix of the
scaled archimedean orbital vectors and the finite last-entry density
then give a bound for $R_{k-1}^2\sum_j\|B_j\|^2$. The next inequality
bounds the real part of the sum over last entries and pairs of heads,
with that density divided by the intermediate group's mass ratio.
Compare \lineSrc{New/Spanning/MainTerm/EnergyEstimate}{422}{archimedean Gram identification},
\lineSrc{New/Spanning/MainTerm/SiegelTopBound}{846}{the fixed-tail Siegel estimate},
\lineSrc{New/Spanning/MainTerm/FiniteSieveSystem}{280}{intermediate group},
\lineSrc{New/Spanning/MainTerm/GlobalSieve}{1145}{global sieve}, and
\lineSrc{New/Spanning/MainTerm/EnergyEstimate}{561}{assembly of the fixed-fixer estimate}.
The manuscript's Appendix F.1 computes the numerical comparison
conditional on the fitted measures and inequalities. The rational
normal form and the error transfer above supply two intermediate
arguments; an independent derivation of the fitted orbital and sieve
bounds is still required.
